% GENERATED FILE. Edit canonical lessons, then run npm run generate:books.
% canonical-source-hash: fnv1a64:f828a5a80f1677c9
% canonical-lessons: RU-C492-otchim, RU-C492-machekha, RU-C492-svodnyy, RU-C492-bliznetsy, RU-C492-dvoyurodnyy, RU-R492-review-485-492

\chapter{Step-parents, Twins and Cousins}
\label{ch:ru-step-parents-twins-and-cousins}
\hlchaptermodality{492}

\begin{chapteropening}
\textbf{By the end of this chapter:} \emph{I can describe a family with step-parents, twins and cousins.}

Complete the last lesson of chapter 492: I can describe a family with step-parents, twins and cousins.
\end{chapteropening}

\section[\texorpdfstring{\ru{отчим} (ótchim)}{отчим (ótchim)}]{\ru{отчим} (ótchim) — a stepfather}

\label{lesson:RU-C492-otchim}



\begin{quote}
Before the new one: say the Russian for a wife's father, then the Russian for a son-in-law.
\end{quote}

\subsection*{You'll want to know: \ru{отчим}}
\textbf{\ru{отчим}} (\emph{ótchim}) — \textquotedblleft{}a stepfather\textquotedblright{}.

\begin{quote}
\ru{Отчим} \ru{относился} \ru{к} \ru{нему} \ru{как} \ru{к} \ru{родному} \ru{сыну}.
\end{quote}

\begin{quote}
\emph{Ótchim otnosílsya k nemú kak k rodnómu sýnu.}
\end{quote}

\begin{quote}
His stepfather treated him like his own son.
\end{quote}

Masculine.

\begin{cousinweb}
\emph{\ru{отчим}} is built on \emph{\ru{отец}} 'father', through the old adjective \emph{\ru{отчий}} 'father's'; EtymDB has no entry for \emph{\ru{отчим}}.
\end{cousinweb}

\begin{grammarlens}[title={What you've built}]
One more word for telling, explaining and describing at B1.
\end{grammarlens}

\subsection*{Your turn}
\begin{itemize}
\raggedright
  \item \emph{Say it:} \emph{ótchim}
  \item \emph{Say it:} \emph{ótchim}, once more
  \item \emph{Say it:} a sentence of your own with \emph{ótchim}
  \item \emph{Recall:} say the Russian for a wife's father, then the Russian for a son-in-law, then say \emph{ótchim} again
\end{itemize}

\begin{tcolorbox}[breakable,colback=teal!4,colframe=teal!35!black,title={Before you move on}]
What does \ru{отчим} mean? (\textquotedblleft{}A stepfather\textquotedblright{}.) Say it once more.
\end{tcolorbox}
\section[\texorpdfstring{\ru{мачеха} (máchekha)}{мачеха (máchekha)}]{\ru{мачеха} (máchekha) — a stepmother}

\label{lesson:RU-C492-machekha}



\begin{quote}
Before the new one: say the Russian for a son-in-law, then the Russian for a stepfather.
\end{quote}

\subsection*{You'll want to know: \ru{мачеха}}
\textbf{\ru{мачеха}} (\emph{máchekha}) — \textquotedblleft{}a stepmother\textquotedblright{}.

\begin{quote}
\ru{В} \ru{этой} \ru{сказке} \ru{мачеха} \ru{была} \ru{очень} \ru{злой}.
\end{quote}

\begin{quote}
\emph{V étoy skázke máchekha bylá óchen zloy.}
\end{quote}

\begin{quote}
In this fairy tale the stepmother was very wicked.
\end{quote}

Feminine.

\begin{cousinweb}
\emph{\ru{мачеха}} is usually explained as \emph{\ru{мать}} 'mother' with the ending \emph{-\ru{еха}}; EtymDB has no entry for \emph{\ru{мачеха}}, and \emph{\ru{мать}} is inherited from Proto-Slavic \emph{mati} (EtymDB).
\end{cousinweb}

\begin{grammarlens}[title={What you've built}]
One more word for telling, explaining and describing at B1.
\end{grammarlens}

\subsection*{Your turn}
\begin{itemize}
\raggedright
  \item \emph{Say it:} \emph{máchekha}
  \item \emph{Say it:} \emph{máchekha}, once more
  \item \emph{Say it:} a sentence of your own with \emph{máchekha}
  \item \emph{Recall:} say the Russian for a son-in-law, then the Russian for a stepfather, then say \emph{máchekha} again
\end{itemize}

\begin{tcolorbox}[breakable,colback=teal!4,colframe=teal!35!black,title={Before you move on}]
What does \ru{мачеха} mean? (\textquotedblleft{}A stepmother\textquotedblright{}.) Say it once more.
\end{tcolorbox}
\section[\texorpdfstring{\ru{сводный} (svódnyy)}{сводный (svódnyy)}]{\ru{сводный} (svódnyy) — step- (as in stepbrother)}

\label{lesson:RU-C492-svodnyy}



\begin{quote}
Before the new one: say the Russian for a stepfather, then the Russian for a stepmother.
\end{quote}

\subsection*{You'll want to know: \ru{сводный}}
\textbf{\ru{сводный}} (\emph{svódnyy}) — \textquotedblleft{}step- (as in stepbrother)\textquotedblright{}.

\begin{quote}
\ru{У} \ru{меня} \ru{есть} \ru{сводный} \ru{брат}.
\end{quote}

\begin{quote}
\emph{U menyá yest svódnyy brat.}
\end{quote}

\begin{quote}
I have a stepbrother.
\end{quote}

\emph{\ru{Сводный} \ru{брат} / \ru{сводная} \ru{сестра}} — a stepbrother, a stepsister.

\begin{cousinweb}
\emph{\ru{сводный}} is built on \emph{\ru{сводить}} 'to bring together', from \emph{\ru{водить}} 'to lead'; EtymDB has no entry for \emph{\ru{сводный}}.
\end{cousinweb}

\begin{grammarlens}[title={What you've built}]
One more word for telling, explaining and describing at B1.
\end{grammarlens}

\subsection*{Your turn}
\begin{itemize}
\raggedright
  \item \emph{Say it:} \emph{svódnyy}
  \item \emph{Say it:} \emph{svódnyy}, once more
  \item \emph{Say it:} a sentence of your own with \emph{svódnyy}
  \item \emph{Recall:} say the Russian for a stepfather, then the Russian for a stepmother, then say \emph{svódnyy} again
\end{itemize}

\begin{tcolorbox}[breakable,colback=teal!4,colframe=teal!35!black,title={Before you move on}]
What does \ru{сводный} mean? (\textquotedblleft{}Step- (as in stepbrother)\textquotedblright{}.) Say it once more.
\end{tcolorbox}
\section[\texorpdfstring{\ru{близнецы} (bliznetsý)}{близнецы (bliznetsý)}]{\ru{близнецы} (bliznetsý) — twins}

\label{lesson:RU-C492-bliznetsy}



\begin{quote}
Before the new one: say the Russian for a stepmother, then the Russian for step-.
\end{quote}

\subsection*{You'll want to know: \ru{близнецы}}
\textbf{\ru{близнецы}} (\emph{bliznetsý}) — \textquotedblleft{}twins\textquotedblright{}.

\begin{quote}
\ru{Близнецы} \ru{так} \ru{похожи}, \ru{что} \ru{их} \ru{путают} \ru{даже} \ru{друзья}.
\end{quote}

\begin{quote}
\emph{Bliznetsý tak pokhózhi, chto ikh pútayut dázhe druzyá.}
\end{quote}

\begin{quote}
The twins are so alike that even their friends mix them up.
\end{quote}

Plural of \emph{\ru{близнец}}; genitive \ru{близнецов}.

\begin{cousinweb}
\emph{\ru{близнецы}} is the plural of \emph{\ru{близнец}}, which is inherited from Proto-Slavic \emph{blizn\ru{ь}c\ru{ь}} (EtymDB) and shares its root with \emph{\ru{близкий}} 'near, close', from Proto-Slavic \emph{bliz\ru{ъ}} 'near' (EtymDB).
\end{cousinweb}

\begin{grammarlens}[title={What you've built}]
One more word for telling, explaining and describing at B1.
\end{grammarlens}

\subsection*{Your turn}
\begin{itemize}
\raggedright
  \item \emph{Say it:} \emph{bliznetsý}
  \item \emph{Say it:} \emph{bliznetsý}, once more
  \item \emph{Say it:} a sentence of your own with \emph{bliznetsý}
  \item \emph{Recall:} say the Russian for a stepmother, then the Russian for step-, then say \emph{bliznetsý} again
\end{itemize}

\begin{tcolorbox}[breakable,colback=teal!4,colframe=teal!35!black,title={Before you move on}]
What does \ru{близнецы} mean? (\textquotedblleft{}Twins\textquotedblright{}.) Say it once more.
\end{tcolorbox}
\section[\texorpdfstring{\ru{двоюродный} (dvoyúrodnyy)}{двоюродный (dvoyúrodnyy)}]{\ru{двоюродный} (dvoyúrodnyy) — related as cousins}

\label{lesson:RU-C492-dvoyurodnyy}



\begin{quote}
Before the new one: say the Russian for step-, then the Russian for twins.
\end{quote}

\subsection*{You'll want to know: \ru{двоюродный}}
\textbf{\ru{двоюродный}} (\emph{dvoyúrodnyy}) — \textquotedblleft{}related as cousins\textquotedblright{}.

\begin{quote}
\ru{Двоюродная} \ru{сестра} \ru{пригласила} \ru{меня} \ru{на} \ru{свадьбу}.
\end{quote}

\begin{quote}
\emph{Dvoyúrodnaya sestrá priglasíla menyá na svádbu.}
\end{quote}

\begin{quote}
My cousin invited me to her wedding.
\end{quote}

\emph{\ru{Двоюродный} \ru{брат} / \ru{двоюродная} \ru{сестра}} — a (first) cousin.

\begin{cousinweb}
\emph{\ru{двоюродный}} is a compound that EtymDB traces to Old East Slavic \emph{d\ru{ъ}voju rodu} 'of two generations', with \emph{\ru{род}} 'kin'.
\end{cousinweb}

\begin{grammarlens}[title={What you've built}]
One more word for telling, explaining and describing at B1.
\end{grammarlens}

\subsection*{Your turn}
\begin{itemize}
\raggedright
  \item \emph{Say it:} \emph{dvoyúrodnyy}
  \item \emph{Say it:} \emph{dvoyúrodnyy}, once more
  \item \emph{Say it:} a sentence of your own with \emph{dvoyúrodnyy}
  \item \emph{Recall:} say the Russian for step-, then the Russian for twins, then say \emph{dvoyúrodnyy} again
\end{itemize}

\begin{tcolorbox}[breakable,colback=teal!4,colframe=teal!35!black,title={Before you move on}]
What does \ru{двоюродный} mean? (\textquotedblleft{}Related as cousins\textquotedblright{}.) Say it once more.
\end{tcolorbox}
\section[(dialogue)]{Words from chapters 485–492 — a review}

\label{lesson:RU-R492-review-485-492}



\begin{quote}
Say the Russian for to feel sick and nausea.
\end{quote}

\subsection*{Your turn}
Say these aloud:

\begin{itemize}
\raggedright
  \item to feel sick, nausea, dizziness, a faint, to have the shivers
  \item a fever, chills, to blow one's nose, to go hoarse, to tickle
  \item a rash, an itch, to itch, a vaccination, to get infected
  \item insomnia, to snore, to yawn, a nightmare, to doze
  \item to examine, a medical examination, a medical test, an X-ray, a surgeon
  \item to recover, recovery, to get better, to swallow, syrup
  \item a husband's mother, a wife's mother, a husband's father, a wife's father, a son-in-law
  \item a stepfather, a stepmother, step-, twins, related as cousins
\end{itemize}

\begin{tcolorbox}[breakable,colback=teal!4,colframe=teal!35!black,title={Before you move on}]
To feel sick? (\textbf{\ru{тошнить}}.) Nausea? (\textbf{\ru{тошнота}}.) Dizziness? (\textbf{\ru{головокружение}}.) A faint? (\textbf{\ru{обморок}}.) To have the shivers? (\textbf{\ru{знобить}}.) A fever? (\textbf{\ru{жар}}.) Chills? (\textbf{\ru{озноб}}.) To blow one's nose? (\textbf{\ru{сморкаться}}.) To go hoarse? (\textbf{\ru{охрипнуть}}.) To tickle? (\textbf{\ru{першить}}.) A rash? (\textbf{\ru{сыпь}}.) An itch? (\textbf{\ru{зуд}}.) To itch? (\textbf{\ru{чесаться}}.) A vaccination? (\textbf{\ru{прививка}}.) To get infected? (\textbf{\ru{заражаться}}.) Insomnia? (\textbf{\ru{бессонница}}.) To snore? (\textbf{\ru{храпеть}}.) To yawn? (\textbf{\ru{зевать}}.) A nightmare? (\textbf{\ru{кошмар}}.) To doze? (\textbf{\ru{дремать}}.) To examine? (\textbf{\ru{обследовать}}.) A medical examination? (\textbf{\ru{обследование}}.) A medical test? (\textbf{\ru{анализ}}.) An X-ray? (\textbf{\ru{рентген}}.) A surgeon? (\textbf{\ru{хирург}}.) To recover? (\textbf{\ru{выздоравливать}}.) Recovery? (\textbf{\ru{выздоровление}}.) To get better? (\textbf{\ru{поправляться}}.) To swallow? (\textbf{\ru{глотать}}.) Syrup? (\textbf{\ru{сироп}}.) A husband's mother? (\textbf{\ru{свекровь}}.) A wife's mother? (\textbf{\ru{тёща}}.) A husband's father? (\textbf{\ru{свёкор}}.) A wife's father? (\textbf{\ru{тесть}}.) A son-in-law? (\textbf{\ru{зять}}.) A stepfather? (\textbf{\ru{отчим}}.) A stepmother? (\textbf{\ru{мачеха}}.) Step-? (\textbf{\ru{сводный}}.) Twins? (\textbf{\ru{близнецы}}.) Related as cousins? (\textbf{\ru{двоюродный}}.)
\end{tcolorbox}
